\scanpage{199}
Meme pour un schema \inserted{local} noetherien, il est possible que la condition
des chaines ne soit pas verifiee, cf Nagata, On the chain problem of
prime ideals, Nagoya Math. Journal, Vol. 10, June 1956, \deleted{\unsure{qui contient|qui traitant}}
(article contenant des resultats interessants sur cette question).

Voir aussi

Nous \deleted{\unsure{verrons}}\inserted{rappelerons} cependant au numero suivant des conditions tres generales
moyennant lesquelles la condition des chaines est satisfaite, conditions
heureusement tou\deleted{\unsure{g}}jours satisfaites en pratique. Quant a l'equidimensionnalite
stricte, elle peut etre en defaut dans des cas non
pathologiques (donc lors meme que la condition des chaines est satisfaite)
pour deux raisons: 1) Les composantes irreductibles de $X$ n'ont
pas meme dimension \deleted{\unsure{et}} (\unsure{on|en} dit alors que $X$ n'est pas \uline{equidimensionnel}\deleted{\unsure{)}},
terminologie qu'il faudrait introduire dans le texte -) 2) Meme si
$X$ est irreductible, et meme si $X$ est un schema integre \uline{regulier}, (i.e.
dont les anneaux locaux sont des anneaux reguliers), il est possible
que les points fermes de $X$ n'aient pas meme codimension. Ainsi soit $V$
un anneau de valuation discrete, soit $K$ son corps des fractions, \inserted{$k$ son corps residuel} $T$ une
indeterminee,posons $A=V\ T$, $X=\mathrm{Spec}(A)$, de sorte que les points
fermes de $X$ correspondent aux ideaux \unsure{maxmaux|maximaux} de $A$, et leur codimension
dans $X$ au \uline{rang} desdits ideaux. Il y a des ideaux maximaux de \deleted{rang}
rang 2, par exemple ceux provenant d'ideaux maximaux de $k\ T$, \deleted{\unsure{ou de}}
\deleted{\uline{certains} ideaux maximaux de \unsure{$K\ T$|$k\ T$}.} Mais \deleted{\unsure{pour}} il y a des ideaux
\unsure{makkmaux|maximaux} de $A$ de rang 1, provenant de\deleted{\unsure{s autres|s}} \inserted{certains} ideaux \unsure{makimaux|maximaux} de $K\ T$,
par exemple l'ideal principal $(aT-1)$, ou $a$ est un generateur de
l'ideal maximal de $V$.

2. \uline{Dimension des preschemas.}

Rappelons la definition suivante:

\uline{Definition} 2. Soient $A$ un anneau commutatif. On appelle dimension de
$A$ la dim. de $\mathrm{Spec}(A)$, et si $J$ est un ideal de $A$, on appelle rang de $J$
la codimension de $V(J)$ dans $\mathrm{Spec}\unsure{!|(}A)$. On dit que $A$ \deleted{est} satisfait
la condition des chianes (resp. est equidim. resp. strictement equidim)
si c'est comme ça pour $\mathrm{Spec}(A)$.

\uline{Proposition 6} \quad Soient $X$ un preschema, $Y$ une partie fermee irreductible
alors on a
\[
\mathrm{codim}_X Y=\dim\underline O_{X/Y}
\]

\uline{Corollaire} 1 \quad La condition des chaines pour $X$ equivaut a la condition
des chaines pour les $\underline O_x$ ($x\unsure{\epsilon|\in}X$); si $X$ est noetherien \deleted{\unsure{on|$X$}}, on peut se
borner aux point fermes.

\uline{Corollaire} 2 \quad Si $X$ est noetherien, on a $\dim X=\mathop{\mathrm{Sup}}\limits_{x\;X}\dim\underline O_x$, et
on peut s'y borner aux point fermes.
% archive sheet number: 198
